%% 第 24 章：等距同构、酉算子与正交算子
%% 本文件由 tools/md2tex.py 自动生成，请勿直接编辑。

\chapter{等距同构、酉算子与正交算子}


\begin{jdefn}{等距映射}{r:25:1}
线性映射：

\[
S:V\to W
\]

若满足：

\[
\|S(v)\|=\|v\|
\]

对所有 $v\in V$ 成立，则称 $S$ 是等距映射。

在线性映射情形，保范等价于保内积：

\[
\langle S(u),S(v)\rangle_W
=
\langle u,v\rangle_V.
\]

这可由实数或复数的极化公式推出。
\end{jdefn}

\begin{jdefn}{等距同构}{r:25:2}
若 $S$ 既是等距映射又是双射，则称 $S$ 是等距同构。

当定义域与陪域是同一内积空间时：

\begin{jitem}
\item %
实数域上的等距同构称为正交算子；
\item %
复数域上的等距同构称为酉算子。
\end{jitem}
\end{jdefn}

\begin{jprop}{等距同构与伴随}{r:25:3}
若：

\[
S:V\to W
\]

是等距同构，则：

\[
S^*=S^{-1}.
\]

因此：

\[
S^*S=I_V,
\qquad
SS^*=I_W.
\]

若 $S:V\to V$，则它是正规算子：

\[
S^*S=SS^*=I.
\]

等距同构一般不自伴。

例如平面上的 $90^\circ$ 旋转是正交算子，但：

\[
S^*\ne S.
\]

只有额外满足：

\[
S=S^*
\]

时，等距同构才同时是自伴的。
\end{jprop}
